\begin{figure}[tbp]
\centering
\begin{tikzpicture}[x=1cm,y=1cm]
\node[box,text width=2.3cm] (va) at (0,0) {guest VA\\取指 / load / store};
\node[box,text width=2.3cm] (tlb) at (3.7,0) {软件 TLB\\检查标签与权限};
\node[greenbox,text width=2.6cm] (fast) at (8.0,0) {RAM 快路径\\由缓存偏移得到 HVA};
\draw[wire] (va)--(tlb); \draw[wire] (tlb)--node[above,font=\footnotesize]{RAM 命中}(fast);
\node[amberbox,text width=2.3cm] (mmu) at (3.7,-1.8) {guest MMU 模型\\查页表与权限};
\node[amberbox,text width=2.3cm] (fault) at (0,-1.8) {guest 访存异常\\交付 guest handler};
\node[box,text width=2.9cm] (as) at (8.0,-1.8) {guest PA\\AddressSpace / MemoryRegion};
\draw[wire] (tlb)--node[left,font=\footnotesize]{miss}(mmu);
\draw[wire] (mmu)--node[above,font=\footnotesize]{失败}(fault);
\draw[wire] (mmu)--node[above,font=\footnotesize]{成功}(as);
\node[greenbox,text width=2.6cm] (ram) at (6.0,-3.7) {RAM 后备内存\\HVA + 页内偏移};
\node[amberbox,text width=2.6cm] (io) at (10.0,-3.7) {MMIO\\调用设备读写回调};
\draw[wire] (as.south)--++(0,-0.45)-|(ram.north);
\draw[wire] (as.south)--++(0,-0.45)-|(io.north);
\node[note] at (4.8,-4.8) {有效映射可填入软件 TLB；宿主 MMU 另行管理 HVA 到 host PA 的转换。};
\end{tikzpicture}
\caption[软件 MMU 的访存路径]{软件 MMU 的访存路径。软件 TLB miss 触发补充转换，不必然产生 guest 页故障；MMIO 和 RAM 共享 guest 物理地址空间，但执行路径不同。\cite{source3}}
\label{fig:softmem}
\end{figure}
